\section*{Реферат}

В рамках выполнения практико-ориентированного проекта распределение задач внутри мини-группы осуществлено следующим образом:
\begin{itemize}
    \item Рахметов А. Р. — выполнение разделов «Организация процесса тестирования», «Разработка тестовых мероприятий» и «Автоматизация тестирования».
    \item Фролов А. А. — выполнение разделов «Тестирование совместимости», «Тестирование адаптивности» и «Применение искусственного интеллекта в тестировании».
    \item Филимонов Е. Е. — выполнение разделов «Введение», «Тестирование верстки», «Технический аудит» и «Заключение».
    \item Исаков А. В. — выполнение разделов «Тестирование сетевых подключений» и «Определение элементов автоматизации тестирования на основе локаторов».
\end{itemize}

\section*{Цель}
Получение практических навыков по тестированию и контролю качества веб-приложения.

\section*{Задачи}
\begin{enumerate}
    \item Разработать план тестирование.
    \item Организовать процесс тестирования с фиксацией задач в системе управления проектами.
    \item Провести тестирование совместимости веб-системы.
    \item Протестировать адаптивность веб-системы.
    \item Проверить валидность верстки.
    \item Провести технический аудит веб-системы.
    \item Проверить качество сетевых подключений.
    \item Автоматизировать процесс выполнения тестовых мероприятий.
    \item Применить нейросетевую модель при проектировании тестовых мероприятий.
    \item Зафиксировать результаты тестирования в отчете
\end{enumerate}

\section*{Предметная область}

В качестве объекта тестирования выбрано веб-приложение Ассоциации \textit{Цифровая Энергетика}. Предметная область данного ресурса охватывает процессы цифровой трансформации электроэнергетики, агрегацию новостей отрасли, предоставление аналитических материалов и координацию экспертных групп. 

URL веб-приложения: \url{https://www.digital-energy.ru/}. Графический интерфейс главной страницы представлен на \cref{fig:main}.

\begin{figure}[H]
    \centering
    \includegraphics[width=\textwidth]{figs/main.png}
    \caption{Главная страница веб-приложения}
    \label{fig:main}
\end{figure}